\section{Continuous potentials and the scalar-boundary action gate}
\label{sec:r24-scalar-boundary-action}

\paragraph{Status and purpose.}
This is a \texttt{FAILED\_ATTEMPT} audit within the frozen RV route.
The action implications proved below are \texttt{DIRECT\_COROLLARY}
consequences of the exact R6 degree-one theorem.  The required
continuous-potential existence statement remains \texttt{OPEN}.
The scalar rank-one interpretation tested below is \texttt{FALSE}
in its stated scope, but it has neither the centered homogeneity
nor the hypotheses of RV15.  No principal theorem or new
counterexample-family count is assigned to this note.  In particular
we do not prove KC12, quadratic Aeppli exactness, or CTO.

Fix the original finite extension \(K/\mathbf Q_p\), its unchanged
absolute value, and an original good, Hausdorff, boundaryless,
quasi-smooth space.  On a sufficiently small ordinary representative
\(V\) assume the dimension is the constant \(n\).
The notation \(B^{p,q}\) denotes the actual Mihatsch delta-form
sheaf, with every actual local refinement retained.  We use the
anticommuting CLD convention \(D=d'd''\), and write \([\alpha]\)
for the ordinary current of an already existing delta-form.
The potential space \(\mathcal D(V)\) is the unchanged space of
finite continuous differences of locally bounded original classical
psh functions.  Every background \(S\) below is symmetric,
middle-positive, and separately \(d'\)- and \(d''\)-closed.
No stronger positive cone or model origin is imposed on \(S\).

\subsection{An ordinary scalar potential already suffices for action}

Write \(M_S(\alpha)\) for the canonical total-degree-one action
of R6, Theorem \texttt{r6-arbitrary-background-degree-one-delta}.
The following observation does not require its potential to define
an actual delta-form.

\begin{proposition}[Continuous-potential action comparison]
Let \(\alpha\in B^{1,1}(V)\) be symmetric.  Suppose there is
\(v\in\mathcal D(V)\) such that
\[
                         [\alpha]=D[v]
             \quad\hbox{as ordinary currents on }V.       \tag{AP24.1}
\]
Then, for every allowed original \(S\) of degree \((p,p)\),
\(0\le p<n\),
\[
                         M_S(\alpha)=D(vS).               \tag{AP24.2}
\]
This is an equality on all original compact smooth complementary
tests.  It holds in every actual presentation and refinement.
Consequently \(M_S(\alpha)\) is separately closed and, whenever
the degree is legal, \(D M_S(\alpha)=0\).
The same assertion is valid for germs: choose the representative
and the potential once, before quantifying over backgrounds and
tests.  No equality on GK delta tests is asserted.
\end{proposition}

\begin{proof}
R6, Lemma \texttt{r6-degree-one-hessian-presentation}, gives
locally a finite presentation
\[
                  \alpha=\sum_{i=1}^{N} f_iD_\delta g_i,
                                                               \tag{AP24.3}
\]
where each \(f_i\) is continuous piecewise smooth and each
\(g_i\) is a continuous local difference of continuous classical
psh functions.  On a compact test support use finitely many such
ordinary neighborhoods and a subordinate original smooth partition.
The actual embedding and continuous multiplication give
\(
 [\alpha]=\sum_i f_iD[g_i]
\)
on each neighborhood.  Subtract (AP24.1).  The resulting ordinary
relation is
\[
                     \sum_i f_iD[g_i]-D[v]=0 .             \tag{AP24.4}
\]
All coefficients in this relation are continuous; the coefficient
of the last curvature is the constant \(-1\).
Expand each local DC potential into its two original psh components.
R6, Theorem \texttt{r6-continuous-curvature-module}, applies to
this exact relation and gives
\[
                     \sum_i f_iD(g_iS)-D(vS)=0 .           \tag{AP24.5}
\]
By the definition and already proved independence of \(M_S\),
the first sum is \(M_S(\alpha)\).  This proves (AP24.2)
locally and on every partitioned test, hence on the whole test.

Both first derivatives of \(D(vS)=d'd''(vS)\) vanish by the
distributional identities \((d')^2=(d'')^2=0\) and
\(d'd''=-d''d'\).  This does not differentiate a nonsmooth
coefficient informally: \(vS\) is a well-defined current by
continuous multiplication with a coefficient-Radon current, and
all derivatives here are distributional.  The R6 statement is
ordinary-local and already independent of all actual refinements,
so no additional test or chart reduction is used.  The germ
assertion follows by choosing a representative for (AP24.1).
\end{proof}

\begin{corollary}[Automatic local DC regularity of a continuous potential]
In the proposition it is enough that \(v\) be finite continuous
and satisfy (AP24.1).  For every point \(x\in V\), one may
choose an ordinary neighborhood \(V_x\subset V\), before
choosing \(S\), on which \(v\in\mathcal D(V_x)\).
The action identity (AP24.2) then holds on all of \(V\) by
ordinary locality.  We do not assert a global DC decomposition
on a previously prescribed \(V\).
\end{corollary}

\begin{proof}
Choose one finite presentation (AP24.3) about \(x\), and shrink
once so that all \(g_i=g_i^+-g_i^-\) have continuous psh
components on that representative.  Choose a relatively compact
ordinary neighborhood \(V_x\) inside it.  The continuous
coefficients are bounded there.  With
\(G=\sum_i(g_i^++g_i^-)\), take a finite \(C\ge0\)
such that \(|f_i|\le C\) on \(V_x\).  Positivity of the
individual curvatures and (AP24.1) give
\[
                -CD[G]\le D[v]\le CD[G]
                    \quad\hbox{on }V_x.                  \tag{AP24.5a}
\]
The functions \(CG+v\) and \(CG-v\) are finite continuous
and have positive original curvature.  The exact frozen R6
lemma \texttt{r6-continuous-positive-classical} makes them
classical psh, locally and hence on \(V_x\).  Therefore
\[
              v=\tfrac12\bigl((CG+v)-(CG-v)\bigr)
                              \in\mathcal D(V_x).          \tag{AP24.5b}
\]
This construction depends on \(\alpha,v,x\), but not on a
background or test.  Apply the proposition on each \(V_x\).
The two sides of (AP24.2) are intrinsic ordinary currents on
\(V\), since multiplication of \(S\) by the given continuous
\(v\) is already defined.  The local identities consequently
glue, proving the last assertion.
\end{proof}

\begin{corollary}[A weaker sufficient target after the RV reduction]
Let \(A\) be the symmetric finitely piecewise-constant PS germ
in RV12, with full \(d'A=0\), and use the accepted PH/QT
preparation.  Thus, on one ordinary representative,
\[
       q\in\mathcal D(V),\qquad
       \gamma\in B^{0,0,1}(V),\qquad
       D_\delta q=A-\gamma .                              \tag{AP24.6}
\]
If \([\gamma]=D[v]\) for some \(v\in\mathcal D(V)\), then
\[
             M_S(A)=D((q+v)S),\qquad D M_S(A)=0            \tag{AP24.7}
\]
for every allowed \(S\), in every legal degree.  In particular
the QA obstruction of \(A\) vanishes for all
\(0\le p\le n-2\).  This conclusion does not require
\(\gamma=0\), a piecewise-linear potential for \(\gamma\),
or an actual delta-form structure on \(v\).
By the preceding corollary it is enough to have a finite
continuous \(v\) satisfying the ordinary current equation,
with local DC decompositions obtained after initial shrinking.
\end{corollary}

\begin{proof}
QT1--QT3 prove (AP24.6), including the full boundary term
\(\gamma\), and prove that the finite piecewise-quadratic
\(q\) is locally an original convex difference.  It follows that
\([A]=D[q+v]\).  Apply the proposition to \(\alpha=A\).
Alternatively, the proposition for \(\gamma\), the frozen
identity \(M_S(D_\delta q)=D(qS)\), and linearity give the
same equality.  The chosen \(V,q,v\) precede every subsequent
choice of \(S\) and compact test.
\end{proof}

Here is the exact quadratic application if the other RV gate has
also been supplied.  For a finite signed PL quadratic polynomial
\(P=\sum_j a_j[h_j][g_j]\), let
\[
 Z=\sum_j a_j E(h_j,g_j),\qquad
 E(h,g)=\tfrac12(d'_Ph\wedge d''_Pg+d'_Pg\wedge d''_Ph).
                                                               \tag{AP24.8}
\]
The frozen Leibniz identity gives
\(
 \widetilde\tau(P)=\tfrac12D_\delta(\sum_j a_jh_jg_j)-Z
\).
QA therefore gives
\(
 \mathcal A_S(P)=-D M_S(Z)
\).
If the exact balanced lift RV4 has been provided, RV5--RV8 replace
\(Z\) by a separately closed \(A\), changing it by full
\(d'\)- and \(d''\)-exact degree-one forms whose QA action is
zero.  Under the additional potential hypothesis in (AP24.6),
(AP24.7) now gives \(\mathcal A_S(P)=0\) for every allowed
\(S\).  This implication is conditional on both existence steps.
It supplies no carrier-CLN certificate: the unchanged IF
action/full-polar distinction prevents that further inference.
Common additional signed mixed tails may be absorbed into \(S\)
by finite positive-tail expansion in legal degrees.  Unrelated
termwise tails are not identified by this observation.

\subsection{The first-derivative relation which is still missing}

For a pure scalar codimension-one \(\gamma\), the R6 module
presentation can be chosen in the form (AP24.3) with continuous
PS \(f_i\) and finite continuous PL \(g_i\).  Every
\(D_\delta g_i\) is fully separately closed.  Thus full
\(d'\gamma=0\) says
\[
               \sum_i d'_Pf_i\wedge D_\delta g_i=0
                         \quad\hbox{in }B^{2,1}(V).        \tag{AP24.9}
\]
For each fixed allowed \(S\), the already proved Radon Leibniz
rule gives the legal ordinary identity
\[
               d'M_S(\gamma)
                    =\sum_i d'_Pf_i\wedge D(g_iS).         \tag{AP24.10}
\]
Indeed each \(D(g_iS)\) is a finite signed expansion of allowed
positive separately closed coefficient-Radon currents, to which
R6's continuous-PS Radon product rule applies.  No boundary of
\(D_\delta g_i\) has been discarded in (AP24.9), and no
derivative of a moving approximation is used in (AP24.10).

The implication from (AP24.9) to vanishing of the right side of
(AP24.10) has not been proved.  R6's scalar-coefficient theorem
applies to \(\sum_i f_iD[g_i]\), not to this covector-coefficient
relation.  We do not define an action on an arbitrary
\(B^{2,1}\) object by writing the right side and assuming
presentation independence.  The ambiguity RV10, where a normal
one-form times a nonzero cap is zero, also prevents unrestricted
coefficient extraction.  If the implication were established for
these particular \(\gamma\), it would already imply their QA
action vanishes, without proving RV15's stronger zero assertion.

\subsection{Why the represented curvature interval does not create a potential}

There is an available order bound, but not the necessary
surjectivity theorem.  On a relatively compact subdomain where
(AP24.3) is fixed, expand \(g_i=g_i^+-g_i^-\) into continuous
psh functions and set \(G=\sum_i(g_i^++g_i^-)\).
Choose \(C\) larger than the absolute values of all coefficients
there.  Positivity of the individual curvatures gives
\[
                       -C D[G]\le[\gamma]\le C D[G].      \tag{AP24.11}
\]
For the RV/QT scalar boundary \(\gamma=A-D_\delta q\), symmetry
and full separate closedness hold.  Hence
\(CD[G]\pm[\gamma]\) are positive separately closed ordinary
currents.  They have not thereby acquired continuous scalar
potentials.

The exact compact set in R6's curvature-interval theorem is
\[
 \mathscr A(G,C;V)
  =\{D[u]:u\in\mathscr P(G,C;V)\},                         \tag{AP24.12}
\]
where \(\mathscr P\) consists of already existing normalized
continuous psh potentials with the specified ambient order bound.
Its weak compactness and closedness show that a weak limit of
members of this represented family remains represented.  Neither
its definition nor its proof identifies it with all positive
closed currents between \(0\) and \(CD[G]\).
To use it here one must first construct a net of represented
curvatures, on the same \(V\) and in a common such interval,
converging on every original test to the required current.
No such net has been constructed.  R6's continuous-positive
converse and HU/HL scalar harmonic recognition start with a
continuous function; they do not supply the missing function
from the current (AP24.11).

\subsection{Intersecting cap branches do not form a rank-one coefficient space}

The following elementary original pressure test limits a possible
interpretation of the desired coefficient sheaf.  It is not a
counterexample to RV15 or to the continuous-potential existence
target.  Work on an ordinary two-torus polyannulus about the joint
radius-one Gauss point \(x\), with
\[
 s_i=\log|T_i|,\qquad
 C_i=D((s_i)_+[U])\quad(i=1,2).                            \tag{AP24.13}
\]
These are genuine positive separately closed original currents,
the ordinary currents of pure scalar codimension-one delta-forms.
Their supports are the complete normal Gauss fibers, by the
frozen CM1/R9 support theorem.  In particular the points
\(\eta_{1,r}\), for \(r>1\) close to one, lie in
\(\operatorname{supp}C_1\) and outside
\(\operatorname{supp}C_2\); the points \(\eta_{r,1}\) have
the reverse property.  Both actual Gauss-radius paths converge
to \(x\).

If \(a,b\in\mathbf R\) with \(a\ne b\), there is no finite
continuous scalar \(c\) on any ordinary representative of this
germ such that
\[
                      aC_1+bC_2=c(C_1+C_2).                \tag{AP24.14}
\]
To prove this, near \(\eta_{1,r}\) the current \(C_2\) is
zero.  Equality would give \((c-a)C_1=0\) there.  A continuous
scalar multiplying a nonzero positive coefficient-Radon current
can vanish only by being zero on its support: otherwise its
nonzero fixed sign on a small neighborhood is detected by a
positive compact test.  Thus \(c(\eta_{1,r})=a\).
Similarly \(c(\eta_{r,1})=b\).  The two limits and continuity
force \(c(x)=a=b\), a contradiction.

Consequently closedness does not turn a multi-branch cap system
into one continuous scalar coefficient relative to its sum, even
before asking for a locally constant coefficient.  Both weights
in this diagnostic are constant and there is no centered
homogeneous vanishing hypothesis.  Also its currents already have
the explicit continuous PL potential
\(a(s_1)_++b(s_2)_+\).  The test only rules out this rank-one
substitution for a branch-sensitive cycle module.  It adds no
KC or RV counterexample.

\subsection{The geometric continuation step is not supplied by tameness alone}

On a fixed full-rank calibrated refinement, a pure scalar
codimension-one presentation of \(\gamma\) can be reduced to
effective scalar coefficients on its cut-cells.  Full
\(d'_P\gamma=0\) makes each such coefficient constant along
that cell.  Artificial subdivision merely restricts this same
coefficient.  A different actual refinement can, however, change
the calibrated branches and their weighted sums.  Across genuine
incidences the closed-cycle condition is a balancing condition;
it is not an assertion that all branch coefficients are one
scalar.  The preceding pressure test explains that distinction
on original analytic data.

The required additional theorem would give a refinement-compatible
cycle-coefficient object for this fixed finite presentation,
with enough ordinary-local specialization toward \(x\) to
force its closed section to vanish when its centered coefficients
vanish there.  We have constructed neither that specialization
nor a finite detector for all later local refining domains.
This is a precise remaining proof step, not a new defined theory.
UR rules out replacing it by a tangent lift of the old Euler
vector on every local cell.  Its counterexample boundary is not
closed and therefore does not decide this coupled continuation
question.

The primary tameness results checked during this attempt have
narrower typed conclusions.  Hrushovski--Loeser's
\emph{Non-archimedean tame topology and stably dominated types},
in the 166-page author version, Theorem 13.4.1, proves local
contractibility of \(B_F(X)\) for a definable subset \(X\) of
an algebraic variety.  This does not identify our absolute
signed-cycle support or its effective coefficient module with
that specified input.  Arbitrary real facet slopes also cannot
simply be inserted into its rational-linear value-group language.

Ducros, \emph{Espaces de Berkovich, polytopes, squelettes et
th\'eorie des mod\`eles}, Theorem 5.1 in the 53-page author
version \texttt{tropicbis.pdf}, constructs skeleta from inverse
images of the \(n\)-torus skeleton under the hypothesis
\(\dim X\le n\).  A complete normal cap with positive-dimensional
analytic fiber is not supplied with such a finite skeleton by
that theorem.  The associated local tropical-image theorem
controls one finite moment image, not its entire original
delta-relation kernel.  The published erratum to the finite
separation argument was also checked; it does not change these
dimension or quotient scopes.

\paragraph{Remaining exact targets and review disclosure.}
The scalar-potential target is: for every RV/QT boundary
\(\gamma\) at an original ordinary germ, find one ordinary
representative and one finite continuous \(v\) with
\([\gamma]=D[v]\) on all original smooth tests.  It is
\texttt{OPEN}; (AP24.5a)--(AP24.5b) would then give local DC
regularity on an initially smaller representative.  The covector
transfer (AP24.9)--(AP24.10) is
another \texttt{OPEN} sufficient action interface.  Neither
would by itself prove the full-polar seminorm conclusion.
The stronger full-actual zero RV15 and the preceding balanced
lift RV4 also remain \texttt{OPEN}.

The reconstruction researcher constructed this note, after root
suggested pursuing an action target and then the branch-continuation
route.  The researcher previously constructed PH; PH, RV, QT,
QA and the R6 action are frozen inputs, with their original review
roles unchanged.  This construction is not an independent review
of itself.  A separate exact-source review is required before
acceptance.
